% !TEX root = ../main.tex
%% Chapitre 14 — Représentations paramétriques et équations cartésiennes
%% Sources : frontend/cours/resume/terminal/resume_representations_parametriques_terminal.txt
%%           frontend/exercices/terminal_spé_Maths/exercice_representations_parametriques_terminal.txt
%% Tous les calculs ont été vérifiés par programme (calcul exact).
%% Exercice 6 : exercice inédit (remplace un doublon de la source).

\begin{ouverture}[parametriques]{14}{Représentations paramétriques}{Partie 2 · Géométrie dans l'espace}
  {Droites\enspace·\enspace Plans\enspace·\enspace Équations cartésiennes\enspace·\enspace Intersections\enspace·\enspace
   Positions relatives}
  \objectif{Déterminer une représentation paramétrique d'une droite et l'utiliser pour tester l'appartenance d'un point.}
  \objectif{Passer d'une représentation paramétrique d'un plan à une équation cartésienne, et inversement.}
  \objectif{Déterminer l'intersection d'une droite et d'un plan, de deux droites, de deux plans.}
  \objectif{Résoudre des problèmes de synthèse : sections, projetés orthogonaux, distances et volumes.}
  \prerequis{Vecteurs, droites et plans de l'espace (chapitre 11).}
  \prerequis{Vecteur normal et équation cartésienne d'un plan (chapitre 13).}
\end{ouverture}
\miseEnPageSouple

%% =====================================================================
%%  L'ESSENTIEL
%% =====================================================================
\essentiel

\essSection{1}{Représentation paramétrique d'une droite}

\noindent
\begin{minipage}[c]{0.58\linewidth}
\begin{theoreme}[Représentation paramétrique]
\raggedright
La droite $d$ passant par $A(x_A\pv y_A\pv z_A)$ et de vecteur directeur $\vec u\,(a\pv b\pv c)$ est l'ensemble des points
$M(x\pv y\pv z)$ tels que $\vect{AM}=t\vec u$, $t\in\R$, c'est-à-dire
\[\begin{cases}x=x_A+ta\\ y=y_A+tb\\ z=z_A+tc\end{cases}\qquad t\in\R.\]
À chaque réel $t$ correspond un unique point de $d$, et réciproquement.
\end{theoreme}
\end{minipage}\hfill
\begin{minipage}[c]{0.38\linewidth}
\centering
\begin{tikzpicture}[scale=1]
  \draw[otNavy,line width=0.9pt] (-1.3,-0.65) -- (3.6,1.8) node[above,font=\small] {$d$};
  \foreach \tt/\lab in {-1/{t=-1},0/{t=0},1/{t=1},2/{t=2}}{%
    \node[otPt] at ({0.4+1.1*\tt},{0.2+0.55*\tt}) {};
    \node[font=\scriptsize,text=otMuted,anchor=north west] at ({0.4+1.1*\tt},{0.2+0.55*\tt}) {$\lab$};}
  \draw[otVec] (0.4,0.2) -- (1.5,0.75) node[midway,above left,font=\small,inner sep=1pt] {$\vec u$};
  \node[otSommet,anchor=south east] at (0.4,0.2) {$A$};
\end{tikzpicture}\\[0.5mm]
\legende{Le paramètre $t$ repère la position du point sur la droite.}
\end{minipage}

\begin{savoirfaire}[Tester l'appartenance d'un point à une droite]
On cherche $t$ à l'aide d'une des trois équations, puis on vérifie que les deux autres sont satisfaites avec la même valeur. Une droite
admet une infinité de représentations paramétriques : on peut changer de point $A$ et remplacer $\vec u$ par tout vecteur non nul qui
lui est colinéaire.
\end{savoirfaire}

\essSection{2}{Représentation paramétrique et équation d'un plan}

\begin{propriete}[Deux descriptions d'un plan]
Le plan passant par $A$ et dirigé par deux vecteurs non colinéaires $\vec u$ et $\vec v$ est l'ensemble des points $M$ tels que
$\vect{AM}=s\vec u+t\vec v$ : c'est une \textbf{représentation paramétrique} à deux paramètres $s$ et $t$. Pour obtenir une
\textbf{équation cartésienne}, on élimine $s$ et $t$, ou l'on cherche un vecteur $\vec n$ orthogonal à $\vec u$ et à $\vec v$
(chapitre 13). Réciproquement, de $ax+by+cz+d=0$ on tire une représentation paramétrique en exprimant une coordonnée en fonction des
deux autres, prises comme paramètres.
\end{propriete}

\essSection{3}{Intersections et positions relatives}

\noindent
\begin{minipage}[c]{0.6\linewidth}
{\renewcommand{\arraystretch}{1.3}\arrayrulecolor{otRule}%
\begin{tabularx}{\linewidth}{@{}>{\raggedright\arraybackslash}p{20mm}Y@{}}
\toprule
\otTete{Objets} & \otTete{Méthode} \\
\midrule
\otNom{Droite et plan} & on reporte la représentation de la droite dans l'équation du plan : une solution $t$ (sécants), aucune
(strictement parallèles) ou une infinité (droite incluse) \\
\otNom{Deux droites} & on résout le système à \textbf{deux paramètres différents} $t$ et $s$ : une solution (sécantes) ; aucune
solution et vecteurs non colinéaires (non coplanaires) \\
\otNom{Deux plans} & normaux colinéaires : parallèles ; sinon sécants suivant une droite, obtenue en prenant une coordonnée comme
paramètre \\
\bottomrule
\end{tabularx}}
\end{minipage}\hfill
\begin{minipage}[c]{0.36\linewidth}
\centering
\begin{tikzpicture}[scale=0.95]
  \fill[otNavyTint] (-0.6,-0.6) -- (3.2,-0.6) -- (4.1,1.3) -- (0.3,1.3) -- cycle;
  \draw[otNavy!45,line width=0.5pt] (-0.6,-0.6) -- (3.2,-0.6) -- (4.1,1.3) -- (0.3,1.3) -- cycle;
  \draw[otNavy,line width=0.9pt,dash pattern=on 2.2pt off 1.6pt] (1.45,-1.25) -- (1.75,0.35);
  \draw[otNavy,line width=0.9pt] (1.75,0.35) -- (2.25,2.4) node[right,font=\small] {$d$};
  \draw[otNavy,line width=0.9pt] (1.33,-1.9) -- (1.45,-1.25);
  \node[otPtV] at (1.75,0.35) {}; \node[otSommet,anchor=north west] at (1.75,0.35) {$I$};
  \node[font=\small,text=otNavy] at (3.6,-0.3) {$\mathcal{P}$};
\end{tikzpicture}\\[0.5mm]
\legende{Une droite sécante à un plan en $I$.}
\end{minipage}

\begin{savoirfaire}[Intersection d'une droite et d'un plan, de deux plans]
\textbf{Droite et plan.} Pour $d:(1+2t\pv-1+t\pv3-t)$ et $\mathcal{P}:x+y+z-5=0$ : $(1+2t)+(-1+t)+(3-t)-5=2t-2=0$, d'où $t=1$ et le
point d'intersection $I(3\pv0\pv2)$.\\[1mm]
\textbf{Deux plans.} $\mathcal{P}_1:x+y+z-3=0$ et $\mathcal{P}_2:x-y+z-1=0$ ont des normaux $(1\pv1\pv1)$ et $(1\pv-1\pv1)$ non colinéaires :
ils sont sécants. Par différence, $2y-2=0$, soit $y=1$, puis $x+z=2$. En posant $z=t$ : la droite d'intersection est
$(2-t\pv1\pv t)$, $t\in\R$, de vecteur directeur $(-1\pv0\pv1)$.
\end{savoirfaire}

\begin{attention}[Pièges fréquents]
\begin{itemize}[label={\color{otOrange}\textbullet},leftmargin=1.4em,itemsep=0.6mm,topsep=0.5mm]
  \item Pour deux droites, il faut deux paramètres de noms différents : un point commun n'a pas le même paramètre sur chaque droite.
  \item Après avoir trouvé les paramètres avec deux équations, on vérifie toujours la troisième.
  \item Une représentation paramétrique n'est pas unique : deux élèves peuvent trouver des systèmes différents et justes.
  \item Dans l'espace, une seule équation $ax+by+cz+d=0$ définit un plan, jamais une droite.
\end{itemize}
\end{attention}

%% =====================================================================
%%  EXERCICES
%% =====================================================================
\exercices
\rubrique{Découverte}

\begin{exercice}[1]{Une droite du plan}
Dans un repère du plan, on considère la droite $d$ passant par $A(3\pv-2)$ et de vecteur directeur $\vec u\,(1\pv2)$.
\begin{questions}
  \item Écrire une représentation paramétrique de $d$.
  \item Vérifier que le point $B(5\pv2)$ appartient à $d$.
  \item Déterminer le point de $d$ de paramètre $t=-2$, puis le point de $d$ d'ordonnée $4$.
  \item Donner un autre vecteur directeur de $d$.
\end{questions}
\end{exercice}

\begin{exercice}[1]{Une droite de l'espace}
On considère la droite $d$ passant par $A(1\pv2\pv-3)$ et de vecteur directeur $\vec u\,(2\pv-1\pv4)$.
\begin{questions}
  \item Écrire une représentation paramétrique de $d$.
  \item Le point $B(3\pv1\pv1)$ appartient-il à $d$ ? Et le point $C(0\pv3\pv-7)$ ?
  \item Déterminer le point de $d$ de paramètre $t=2$, puis le point de $d$ de cote $z=5$.
\end{questions}
\end{exercice}

\begin{exercice}[1]{Plusieurs représentations d'une même droite}
On considère la droite $d$ passant par $A(2\pv-1\pv3)$ et de vecteur directeur $\vec u\,(1\pv2\pv-1)$.
\begin{questions}
  \item Donner une représentation paramétrique de $d$.
  \item Les points $M(4\pv3\pv1)$ et $N(3\pv1\pv2)$ appartiennent-ils à $d$ ? Que représente $\vect{AM}$ pour la droite ?
  \item Déterminer le point $P$ de $d$ de paramètre $t=-2$.
  \item Donner une autre représentation paramétrique de $d$, utilisant le point $M$ et le vecteur $\vect{AM}$. Quel paramètre
    correspond alors au point $A$ ?
\end{questions}
\end{exercice}

\rubrique{Entraînement}

\begin{exercice}[2]{Du paramétrique au cartésien dans le plan}
Dans un repère orthonormé du plan, on considère la droite $d:\begin{cases}x=1+2t\\y=3-4t\end{cases}$, $t\in\R$.
\begin{questions}
  \item En exprimant $t$ en fonction de $x$, déterminer une équation cartésienne de $d$. En déduire un vecteur normal à $d$.
  \item Déterminer une équation de la droite $d'$ parallèle à $d$ passant par l'origine $O$.
  \item Le projeté orthogonal $H$ de $O$ sur $d$ est sur la droite dirigée par le vecteur normal et passant par $O$. Déterminer $H$ et
    en déduire la distance entre les droites parallèles $d$ et $d'$.
\end{questions}
\end{exercice}

\begin{exercice}[2]{Équation cartésienne d'un plan}
L'espace est muni d'un repère orthonormé. On considère le plan $\mathcal{P}$ passant par $A(1\pv2\pv-1)$ et de vecteur normal
$\vec n\,(3\pv-1\pv2)$.
\begin{questions}
  \item Déterminer une équation cartésienne de $\mathcal{P}$. Les points $B(2\pv5\pv1)$ et $O$ appartiennent-ils à $\mathcal{P}$ ?
  \item Vérifier que $\vec{u_1}\,(1\pv3\pv0)$ et $\vec{u_2}\,(0\pv2\pv1)$ dirigent $\mathcal{P}$. En déduire une représentation
    paramétrique de $\mathcal{P}$.
  \item Déterminer une équation du plan parallèle à $\mathcal{P}$ passant par l'origine.
\end{questions}
\end{exercice}

\begin{exercice}[2]{D'une représentation à l'autre}
L'espace est muni d'un repère orthonormé.
\begin{questions}
  \item Le plan $\mathcal{P}$ passe par $A(1\pv1\pv0)$ et il est dirigé par $\vec u\,(1\pv0\pv2)$ et $\vec v\,(0\pv1\pv-1)$. Écrire une
    représentation paramétrique de $\mathcal{P}$, puis, en éliminant les paramètres, une équation cartésienne.
  \item Vérifier que le vecteur $\vec n\,(2\pv-1\pv-1)$ est orthogonal à $\vec u$ et à $\vec v$. Est-ce cohérent ?
  \item Le point $M(3\pv2\pv3)$ appartient-il à $\mathcal{P}$ ? Si oui, déterminer ses paramètres.
  \item Donner une représentation paramétrique du plan $\mathcal{Q}:x+y+z=3$.
\end{questions}
\end{exercice}

\begin{exercice}[2]{Quatre droites et leurs positions}
L'espace est muni d'un repère. On considère :
\begin{itemize}[leftmargin=1.4em,itemsep=0.3mm,topsep=0.5mm]
  \item la droite $d_1$ passant par $A(1\pv-1\pv2)$ et de vecteur directeur $\vec u\,(2\pv1\pv-1)$ ;
  \item la droite $d_2$ passant par $B(0\pv2\pv1)$ et $C(3\pv-1\pv4)$ ;
  \item la droite $d_3:\begin{cases}x=3+s\\y=-2s\\z=1+s\end{cases}$, $s\in\R$ ;
  \item la droite $d_4$ passant par $F(-1\pv0\pv3)$ et parallèle à $d_1$.
\end{itemize}
\begin{questions}
  \item Déterminer une représentation paramétrique de $d_1$, $d_2$ et $d_4$, ainsi qu'un point et un vecteur directeur de $d_3$.
  \item Les droites $d_1$ et $d_2$ sont-elles coplanaires ?
  \item Montrer que $d_1$ et $d_3$ sont sécantes et déterminer leur point d'intersection $E$.
  \item Montrer que $d_2$ et $d_4$ sont sécantes et déterminer leur point d'intersection $G$.
\end{questions}
\end{exercice}

\rubrique{Problèmes}

\begin{exercice}[2]{Une droite et un plan}
L'espace est muni d'un repère orthonormé. On considère la droite $d:\begin{cases}x=1+2t\\y=3-t\\z=-1+t\end{cases}$, $t\in\R$, et le plan
$\mathcal{P}:x+2y-z+4=0$.
\begin{questions}
  \item Donner un vecteur directeur $\vec u$ de $d$ et un vecteur normal $\vec n$ de $\mathcal{P}$. Calculer $\vec u\cdot\vec n$.
  \item En déduire que $d$ et $\mathcal{P}$ sont sécants, puis déterminer leur point d'intersection $I$.
  \item Déterminer le projeté orthogonal $H$ du point $A(1\pv3\pv-1)$ sur $\mathcal{P}$, puis la distance de $A$ au plan $\mathcal{P}$.
\end{questions}
\end{exercice}

\begin{exercice}[2]{Deux droites sécantes}
L'espace est muni d'un repère orthonormé. On considère $d_1:\begin{cases}x=t\\y=2t\\z=1+t\end{cases}$ et
$d_2:\begin{cases}x=1+s\\y=2+s\\z=2+s\end{cases}$, avec $t,\ s\in\R$.
\begin{questions}
  \item Déterminer un vecteur directeur de chaque droite. Les droites sont-elles parallèles ?
  \item Montrer que $d_1$ et $d_2$ sont sécantes et déterminer leur point d'intersection $I$.
  \item Calculer, au degré près, l'angle formé par les vecteurs directeurs des deux droites.
\end{questions}
\end{exercice}

\begin{exercice}[2]{Deux droites non coplanaires}
On considère $d_1:\begin{cases}x=1+2t\\y=3-t\\z=-1+t\end{cases}$ et $d_2:\begin{cases}x=2+s\\y=1+2s\\z=3s\end{cases}$, avec
$t,\ s\in\R$.
\begin{questions}
  \item Les droites $d_1$ et $d_2$ sont-elles parallèles ?
  \item Résoudre le système formé par les deux premières équations, puis tester la troisième.
  \item Conclure sur la position relative de $d_1$ et $d_2$. Peut-on trouver un plan qui les contienne toutes les deux ?
\end{questions}
\end{exercice}

\begin{exercice}[3]{Deux plans sécants}
L'espace est muni d'un repère orthonormé. On considère les plans $\mathcal{P}_1:x+y+z-1=0$ et $\mathcal{P}_2:2x-y+z+3=0$.
\begin{questions}
  \item Les plans $\mathcal{P}_1$ et $\mathcal{P}_2$ sont-ils parallèles ? perpendiculaires ?
  \item En prenant $y=t$ comme paramètre, déterminer une représentation paramétrique de leur droite d'intersection $\Delta$.
  \item Vérifier que le point $M(-4\pv0\pv5)$ appartient aux deux plans, et que le vecteur directeur de $\Delta$ est orthogonal aux
    deux vecteurs normaux.
  \item Calculer, au degré près, l'angle formé par les vecteurs normaux des deux plans.
\end{questions}
\end{exercice}

\begin{exercice}[3]{Synthèse : droites, plan et intersection}
L'espace est muni d'un repère orthonormé. On considère les points $A(1\pv0\pv2)$, $B(3\pv1\pv0)$, $C(2\pv-1\pv3)$ et la droite
$d:\begin{cases}x=1+t\\y=2-t\\z=3+2t\end{cases}$, $t\in\R$.
\begin{questions}
  \item Déterminer une représentation paramétrique de la droite $(AB)$.
  \item Montrer que $\vec n\,(1\pv4\pv3)$ est normal au plan $(ABC)$, puis que ce plan a pour équation $x+4y+3z-7=0$.
  \item Le point $D(0\pv-1\pv4)$ appartient-il au plan $(ABC)$ ?
  \item Montrer que $d$ coupe le plan $(ABC)$ et déterminer les coordonnées du point d'intersection $I$.
\end{questions}
\end{exercice}

\begin{exercice}[3]{Une pyramide à base rectangulaire}
\begin{codecote}[0.56]
Dans un repère orthonormé, on considère la pyramide $SABCD$ avec $A(0\pv0\pv0)$, $B(3\pv0\pv0)$, $C(3\pv2\pv0)$, $D(0\pv2\pv0)$ et
$S(1\pv1\pv3)$.
\begin{questions}
  \item Montrer que $ABCD$ est un rectangle et calculer son aire.
  \item Déterminer une représentation paramétrique des droites $(SD)$ et $(SA)$.
  \item $P$ est le point de $[SD]$ tel que $\vect{SP}=\frac14\vect{SD}$. Calculer ses coordonnées.
  \item Vérifier que $P$, $B$, $C$ ne sont pas alignés et que le plan $(PBC)$ a pour équation $x+z-3=0$.
\end{questions}
\cote
\centering
\begin{tikzpicture}[otEspace,y={(0.55cm,0.42cm)},scale=1.05]
  \coordinate (A) at (0,0,0); \coordinate (B) at (3,0,0); \coordinate (C) at (3,2,0); \coordinate (D) at (0,2,0);
  \coordinate (S) at (1,1,3);
  \coordinate (P) at (0.75,1.25,2.25); \coordinate (Q) at (0.75,0.75,2.25);
  \fill[otGreen,opacity=0.2] (B) -- (C) -- (P) -- (Q) -- cycle;
  \draw[otVu] (A) -- (B) -- (C) -- (S) -- cycle; \draw[otVu] (S) -- (B);
  \draw[otCache] (A) -- (D) -- (C); \draw[otCache] (D) -- (S);
  \draw[otGreenDark,line width=0.8pt] (B) -- (Q);
  \draw[otGreenDark,line width=0.8pt,dash pattern=on 2.2pt off 1.6pt] (Q) -- (P) -- (C);
  \foreach \p/\a in {A/north east,B/north west,C/west,D/south east,S/south}{\node[otSommet,anchor=\a] at (\p) {$\p$};}
  \node[otPtV] at (P) {}; \node[otSommet,anchor=south east] at (P) {$P$};
  \node[otPtV] at (Q) {}; \node[otSommet,anchor=east] at (Q) {$Q$};
\end{tikzpicture}
\end{codecote}
\begin{questions}[start=5]
  \item Montrer que la droite $(SA)$ coupe le plan $(PBC)$ en un point $Q$ dont on donnera les coordonnées.
  \item Déterminer une représentation paramétrique de la droite $\Delta$ orthogonale au plan $(PBC)$ passant par $S$, puis les
    coordonnées du point $H$, intersection de $\Delta$ et de $(PBC)$. En déduire la distance du point $S$ au plan $(PBC)$.
  \item Le projeté orthogonal de $P$ sur la droite $(BC)$ est le point $K\left(3\pv\frac54\pv0\right)$. Calculer $PK$, puis l'aire du
    triangle $PBC$, et en déduire le volume du tétraèdre $SPBC$.
\end{questions}
\end{exercice}

\begin{exercice}[3]{Un tétraèdre trirectangle}
L'espace est muni d'un repère orthonormé d'origine $O$. On considère les points $A(1\pv0\pv0)$, $B(0\pv2\pv0)$ et $C(0\pv0\pv3)$ :
les arêtes $[OA]$, $[OB]$ et $[OC]$ du tétraèdre $OABC$ sont deux à deux perpendiculaires.
\begin{questions}
  \item Vérifier que le plan $(ABC)$ a pour équation $6x+3y+2z-6=0$.
  \item Déterminer une représentation paramétrique de la droite $\Delta$ passant par $O$ et orthogonale au plan $(ABC)$, puis les
    coordonnées du point $H$, intersection de $\Delta$ et de $(ABC)$.
  \item En déduire la distance du point $O$ au plan $(ABC)$.
  \item Calculer le volume du tétraèdre $OABC$, puis en déduire l'aire du triangle $ABC$.
  \item Montrer que $\vect{AH}\cdot\vect{BC}=0$ et $\vect{BH}\cdot\vect{AC}=0$. Que représente le point $H$ pour le triangle $ABC$ ?
\end{questions}
\end{exercice}
